\section{A divided-difference identity across depths and colors}
\label{logres:sec:main}
The explicit logistic density suggests an identity search that keeps the
depth parameter symbolic instead of recognizing one numerical value at a
time. Its beta moments already determine the one-resolvent transform.
Partial fractions then introduce several independent colors, while
coalescing those colors produces finite formulas at every resolvent order.
The resulting Gaussian logarithmic moments are exact polynomials in
$\pi$ and $\log2$, with no integer-relation premise.

For $-1<\Real t<1$ put
\[
 W_t(u)=\left(\frac{u}{1-u}\right)^t\frac{\sin\pi t}{\pi t},
 \qquad 0<u<1,
\]
using the removable value at $t=0$. This is the generating density
\eqref{ado:eq:kernelgenerator}. Its integral is one, even for complex $t$;
for real $-1<t<1$ it is positive. Denote divided differences by
$[z_1,\ldots,z_m]f$, with the continuous confluent interpretation when
nodes coincide.

\begin{theorem}[Several resolvents and their confluent limits]
\label{logres:thm:colors}
Let $z_1,\ldots,z_m\in\Omega$, where $\Omega=\C\setminus[1,\infty)$.
For distinct nodes and $m\ge2$,
\begin{equation}\label{logres:eq:colors}
 \int_0^1\frac{W_t(u)}{\prod_{j=1}^m(1-z_ju)}\,du
 =\frac1t[z_1,\ldots,z_m]
       \{z^{m-2}((1-z)^{-t}-1)\}.
\end{equation}
For $m=1$ the transform is
\begin{equation}\label{logres:eq:one}
 \int_0^1\frac{W_t(u)}{1-zu}\,du
       =\frac{(1-z)^{-t}-1}{tz}.
\end{equation}
All powers use the principal logarithm; the apparent singularities at
$t=0$, $z=0$, and coincident nodes are removable. In particular, for
$m\ge2$,
\begin{align}\label{logres:eq:confluent}
 \int_0^1\frac{W_t(u)}{(1-zu)^m}\,du
 &=(1-z)^{-t-m+1}\mathcal P_m(t,z),\notag\\
 \mathcal P_m(t,z)
 &=\sum_{j=0}^{m-2}\binom{m-2}{j}
       \frac{(t+1)_j}{(j+1)!}
       z^j(1-z)^{m-2-j},
\end{align}
where $(t+1)_j$ is the rising factorial and $(t+1)_0=1$.
\end{theorem}
\begin{proof}
The beta identity \eqref{ado:eq:mellin} gives
\[
 \int_0^1u^nW_t(u)\,du=\frac{(1+t)_n}{(n+1)!}.
\]
For $|z|<1$, expand the geometric kernel and sum the binomial series.
Its result is \eqref{logres:eq:one}. The integrable endpoint majorants
$u^{-1+\epsilon}(1-u)^{-1+\epsilon}$, on smaller parameter strips,
justify local holomorphy and continuation to the principal slit plane.
For distinct nonzero nodes, the exact partial fractions are
\[
 \frac1{\prod_j(1-z_ju)}
 =\sum_{j=1}^m
   \frac{z_j^{m-1}}{\prod_{k\ne j}(z_j-z_k)}\frac1{1-z_ju}.
\]
Integrating and using \eqref{logres:eq:one} gives the Lagrange formula
for \eqref{logres:eq:colors}. Continuity allows zero nodes.

When every node tends to $z$, the divided difference becomes the
$(m-1)$st derivative divided by $(m-1)!$. Apply the product rule to
$z^{m-2}((1-z)^{-t}-1)/t$. The term differentiating the second factor
zero times vanishes, because the polynomial has degree $m-2$.
For $j+1\ge1$ derivatives of the second factor one has
\[
 \frac{d^{j+1}}{dz^{j+1}}
       \frac{(1-z)^{-t}-1}{t}
       =(t+1)_j(1-z)^{-t-j-1}.
\]
Collecting the binomial factors proves \eqref{logres:eq:confluent}.
All confluent limits also follow directly from domination of the
original integral on compact subsets of $\Omega$.
\end{proof}

For example, the first three confluent cases are
\begin{align*}
 \mathcal P_2&=1,\\
 \mathcal P_3&=1+(t-1)z/2,\\
 \mathcal P_4&=(1-z)^2+(t+1)z(1-z)
                         +(t+1)(t+2)z^2/6.
\end{align*}
Thus the squared-resolvent identity has the particularly simple form
\[
 \int_0^1\frac{W_t(u)}{(1-zu)^2}\,du=(1-z)^{-t-1}.
\]
The several-color identity is symmetric in its nodes; repeated colors
are derivatives, rather than a separate guessed formula.

\begin{corollary}[Every Gaussian logarithmic moment]
\label{logres:cor:Gaussian}
For integers $m\ge2$ and $k\ge0$, define
\[
 I_{m,k}=\int_0^1\frac{\log^k(u/(1-u))}{(1-\iu u)^m}\,du,
 \qquad \lambda=\tfrac12\log2-\tfrac{\iu\pi}4.
\]
These are ordinary absolutely convergent integrals, and
\begin{equation}\label{logres:eq:Gaussian}
 I_{m,k}=\frac{k!}{(1-\iu)^{m-1}}
 [t^k]\left\{\mathcal P_m(t,\iu)
   \exp\left(-\lambda t+
       \sum_{j\ge1}\frac{\zeta(2j)}j t^{2j}\right)\right\}.
\end{equation}
In particular $I_{m,k}\in\Q[\iu,\pi,\log2]$ for every $m,k$.
For the squared resolvent,
\begin{align*}
 I_{2,0}&=\frac1{1-\iu},&
 I_{2,1}&=-\frac{\lambda}{1-\iu},\\
 I_{2,2}&=\frac{\lambda^2+\pi^2/3}{1-\iu},&
 I_{2,3}&=-\frac{\lambda^3+\pi^2\lambda}{1-\iu},\\
 I_{2,4}&=\frac{\lambda^4+2\pi^2\lambda^2+7\pi^4/15}{1-\iu}.
\end{align*}
\end{corollary}
\begin{proof}
Multiply \eqref{logres:eq:confluent} by $\pi t/\sin\pi t$ and
differentiate $k$ times at zero. Powers of the endpoint logarithms
are integrable, so this differentiation is legitimate. The reflection
formula and its local logarithmic expansion give
\[
 \frac{\pi t}{\sin\pi t}=\Gamma(1+t)\Gamma(1-t)
    =\exp\left(\sum_{j\ge1}\frac{\zeta(2j)}j t^{2j}\right),
 \qquad |t|<1.
\]
At $z=\iu$, the principal logarithm is
$\log(1-\iu)=\lambda$. This proves the coefficient formula and the
listed values. Only finitely many terms enter a coefficient of order
$k$, and every even zeta value is a rational multiple of $\pi^{2j}$.
\end{proof}


For a concrete weight-two pair, the real and imaginary parts of
$I_{2,2}$ give the ordinary integral identities
\begin{align}\label{logres:eq:real-second-moments}
 \int_0^1\frac{(1-u^2)\log^2(u/(1-u))}{(1+u^2)^2}\,du
 &=\frac{\log^2 2}{8}+\frac{13\pi^2}{96}
                            +\frac{\pi\log2}{8},\notag\\
 \int_0^1\frac{2u\log^2(u/(1-u))}{(1+u^2)^2}\,du
 &=\frac{\log^2 2}{8}+\frac{13\pi^2}{96}
                            -\frac{\pi\log2}{8}.
\end{align}
Both follow by expanding $\lambda^2$ in the proved formula, so their
constants are derived rather than recognized numerically.

\begin{corollary}[Two real Gaussian generating identities]
\label{logres:cor:real-generators}
For real $-1<t<1$,
\begin{align}\label{logres:eq:real-generators}
 \int_0^1\frac{(u/(1-u))^t}{1+u^2}\,du
 &=\frac{\pi\,2^{-t/2}\sin(\pi t/4)}{\sin\pi t},\notag\\
 \int_0^1\frac{u(u/(1-u))^t}{1+u^2}\,du
 &=\frac{\pi\{1-2^{-t/2}\cos(\pi t/4)\}}{\sin\pi t}.
\end{align}
At $t=0$ the removable values are respectively $\pi/4$ and
$\log2/2$. Differentiating at zero supplies every real logarithmic
moment of these two kernels.
\end{corollary}
\begin{proof}
Use \eqref{logres:eq:one} at $z=\iu$ and $z=-\iu$ and the identities
\[
 \frac1{1+u^2}=\frac12\left(\frac1{1-\iu u}+\frac1{1+\iu u}\right),
 \quad
 \frac{u}{1+u^2}=\frac1{2\iu}\left(\frac1{1-\iu u}-\frac1{1+\iu u}\right).
\]
Multiply by $\pi t/\sin\pi t$ and simplify the conjugate principal
powers. Continuity gives the zero-parameter values.
\end{proof}

The finite confluent polynomials and moment recurrence are checked exactly
in \src{verification/certify_logistic_resolvents.py}. Independent
logistic-coordinate quadratures are labeled as numerical diagnostics;
the arbitrary-order identities follow from the beta moments, partial
fractions and analytic parameter differentiation. These formulas add a
several-color integral family to the development; they do not settle the
separate mixed harmonic sums $S_6$ or $S_8$, and no literature-priority
claim is made for this classical beta-integral mechanism.
